% Folder 83, batch 12 (archivists' page 221, a single page).
% Pass: Opus 5.5 (claude-opus-5-5), 2026-10-10 — first pass, unchecked against the pages by a human.
%
% The one page of this batch is not mathematics: it is a page of a typescript
% of Grothendieck's own, numbered « L 19 » at the head, bearing a few
% corrections in blue ink in his hand. It is a passage of reflection written
% around Récoltes et Semailles (« l'écriture des Quatre Opérations »): the
% « dégradation » of the mathematical milieu since his departure, the contempt
% for the creative, « feminine » side of discovery (dream, vision, idea), the
% « aplatissement » of mathematical thought (Siegel's word). Kept under the
% rule that a typed leaf carrying his hand is transcribed. Typewriter
% x-outs are given as \struck; his blue-ink insertions as \add.

\documentclass[11pt,a4paper]{article}
\input{../preamble/grothendieck}

\folder{83}
\batch{12}
\pages{221}{221}
\foldertitle{[Icosaèdres, bi-icosaèdres et algèbres d'Azumaya de rang 4] : notes manuscrites (1982-1983, 1986, s.d.).}
\dating{1982-1986}
\watermark{Édition de démonstration}

\begin{document}

\page{221}
\note{feuillet dactylographié de sa main, paginé « L 19 » en tête, avec quelques corrections manuscrites à l'encre bleue ; ce n'est pas une page de mathématiques. Le texte commence au milieu d'une phrase venue de la page précédente du dactylogramme, absente de ce lot, et s'interrompt en fin de page au milieu d'une parenthèse. Les mots dactylographiés en lettres espacées sont rendus en italique ; les renvois (*), (**), … visent des notes qui ne figurent pas sur ce feuillet.}

Récoltes et Semailles j'ai pu me rendre compte, et de façon saisissante en ces tout derniers mois surtout (avec l'écriture des Quatre Opérations), qu'il y a eu depuis mon départ de la scène mathématique une stupéfiante \emph{dégradation} dans l'esprit qui aujourd'hui fait loi dans les milieux que j'avais connus, et (me semble-t-il, dans une large mesure au moins) dans le monde mathématique en général (*). Il est possible même, \struck{\ill{}} autant par ma personnalité mathématique très particulière que par les conditions qui ont entouré mon départ, que celui-ci ait agi comme un catalyseur dans une évolution qui était déjà en train de se faire (**) -- une évolution dont je n'ai \add{alors} rien su percevoir (pas plus qu'aucun autre de mes collègues et amis, à la seule exception peut-être de Claude Chevalley).
L'aspect de cette dégradation auquel je pense surtout ici (qui \add{en} est \add{juste} un aspect parmi de nombreux autres (***)) est le \emph{mépris tacite}, quand ce n'est la dérision sans équivoque, à l'encontre de ce qui (en mathématique, en l'occurence) ne s'apparente pas au pur travail du marteau sur l'enclume ou sur le burin -- le mépris des processus créateurs les plus délicats (et \struck{parf} souvent de moindre apparence\add{)} ; de tout ce qui est \struck{inspiration,} \emph{inspiration}, \emph{rêve}, \emph{vision} (si puissantes et si fertiles soient-elles), et même (à la limite) de toute \emph{idée}, si clairement conçue et formulée soit-elle : de tout ce qui n'est écrit et \emph{publié} noir sur blanc, sous forme d'énoncés purs et durs, répertoriables et répertoriés, mûrs pour les « banques de données » \uncertain{engouffrés} dans les inépuisables mémoires de nos mégaordinateurs.\note{le mot « engouffrés » est frappé par-dessus une première frappe que les surcharges rendent illisible.}

Il y a eu (pour reprendre une expression de C.L. Siegel (****)) un extraordinaire \struck{\ill{}} « \emph{aplatissement} », un « \emph{rétrécissement} » de la pensée mathématique, dépouillée d'une dimension essentielle, de tout son « versant d'ombre », du versant « féminin ». Il est vrai que par une tradition ancestrale, ce versant-là du travail de découverte restait dans une large mesure occulté, personne (autant dire) n'en \emph{parlait} jamais -- mais le contact \add{vivant} avec les sources profondes du rêve, qui alimentent les grandes visions et les grands desseins, n'avait jamais encore (à ma connaissance) été perdu. Il semblerait que \add{dès à} présent nous soyons déjà entrés dans une \emph{époque de dessèchement}, où cette source est, non point tarie certes, mais où l'accès à elle est condamné\struck{e}, par le \emph{verdict sans} appel du mépris général et par les représailles de la \emph{dérision}.\note{« verdict sans » et « dérision » sont soulignés à l'encre bleue ; deux traits obliques à l'encre bleue, après « aplatissement » et entre « mesure » et « occulté », séparent des mots ; un double trait vertical bleu marque ensuite un alinéa.}

Nous voilà approcher du moment, semble-t-il, où sera éradiqué en chacun non seulement le \emph{souvenir} de tout travail proche de la source, du travail « au féminin » (ridiculisé
\end{document}
